\documentclass[a4paper,12pt]{article}
\usepackage[T2A]{fontenc}
\usepackage[utf8]{inputenc}
\usepackage[russian]{babel}
\usepackage{amsmath,amssymb,mathtools}
\usepackage[margin=22mm,headheight=15pt,headsep=9mm,footskip=14mm]{geometry}
\usepackage{microtype}
\usepackage{xcolor}
\usepackage{booktabs,tabularx,array}
\usepackage{hyperref}
\usepackage{tikz}
\usepackage{fancyhdr}
\usepackage{parskip}
\definecolor{accent}{HTML}{008080}
\definecolor{darkaccent}{HTML}{004D4D}
\definecolor{bodygray}{HTML}{3E4648}
\hypersetup{colorlinks=true,linkcolor=darkaccent,pdfauthor={Лёвин Артём Александрович},pdftitle={Ревью домашней работы — Надежда Клименко — 20.11.2026}}
\setlength{\parindent}{0pt}
\setlength{\parskip}{6pt}
\renewcommand{\arraystretch}{1.35}
\pagestyle{fancy}
\fancyhf{}
\fancyhead[L]{\footnotesize\color{bodygray}Отношения, доли и проценты}
\fancyhead[R]{\footnotesize\color{bodygray}Надежда Клименко}
\fancyfoot[C]{\footnotesize\color{bodygray}\thepage}
\renewcommand{\headrulewidth}{0.3pt}
\renewcommand{\footrulewidth}{0pt}
\newcommand{\point}[1]{\par\vspace{2pt}{\bfseries\color{darkaccent}#1}\par}
\begin{document}

\thispagestyle{empty}
\begin{center}
\vspace*{22mm}
{\large\color{darkaccent}МАТЕМАТИКА}\par
\vspace{23mm}
{\LARGE Ревью домашней работы}\par
\vspace{10mm}
{\Large\color{accent}Отношения, доли и проценты}\par
\vspace{9mm}
{\normalsize 6 класс}\par
\vfill
{\large Надежда Клименко}\par
\vspace{13mm}
{\normalsize 20 ноября 2026 г.}\par
\vspace{18mm}
{\small Лёвин Артём Александрович}\par
{\small эксклюзивно для Надежды Клименко}\par
\vspace{20mm}
\begin{tikzpicture}
\draw[accent,line width=0.85pt] (0,0) .. controls (1.2,0.50) and (2.2,-0.5) .. (3.45,0)
  .. controls (4.6,0.50) and (5.45,-0.50) .. (6.6,0)
  .. controls (7.5,0.35) and (8.8,-0.36) .. (9.8,0);
\end{tikzpicture}
\vspace*{6mm}
\end{center}
\clearpage

\section*{Сводная таблица}
\noindent\textcolor{bodygray}{В таблице перечислены только задания, требующие исправления или завершения. Номер задания ведёт к подробному разбору.}

\begingroup
\small\setlength{\tabcolsep}{4pt}
\begin{tabularx}{\linewidth}{@{}p{13mm}p{25mm}X p{22mm}p{29mm}@{}}
\toprule
\textbf{№} & \textbf{Тема} & \textbf{Суть ошибки} & \textbf{Ваш ответ} & \textbf{Правильный ответ} \\
\midrule
\hyperref[task:9]{9} & Процентный рост & Ошибка при вычислении результата по правильно составленной формуле & $60000\%$ & $25\%$ \\
\addlinespace[3pt]
\hyperref[task:10]{10} & Сравнение долей & Решение и ответ отсутствуют & Не указан & $40{,}8\%$ \\
\bottomrule
\end{tabularx}
\endgroup

\vspace{13mm}
\textcolor{darkaccent}{\bfseries Главный ориентир}\par
Когда требуется найти, на сколько процентов одна величина изменилась \emph{относительно другой}, нужно разделить изменение на исходную величину. Когда требуется узнать, сколько процентов одна группа составляет \emph{от другой группы}, в знаменателе должно стоять количество людей во второй группе.
\clearpage

\thispagestyle{empty}
\vspace*{0.34\textheight}
\begin{center}
{\LARGE\color{darkaccent}Разбор ошибок}\par
\vspace{7mm}
{\normalsize Задания № 9 и № 10}
\end{center}
\clearpage

\phantomsection\label{task:9}
\section*{Задание № 9. Процентное увеличение}
\point{Текст задания}
В прошлом году собрали 480 кг макулатуры, в этом~--- 600 кг. На сколько процентов увеличился сбор?

\point{Ваше решение}
В записи верно составлено выражение
\[
  \frac{600-480}{480}\cdot100\%,
\]
после которого записан ответ $60000\%$.

\point{Ваш ответ}
$60000\%$.

\point{Ошибка}
\textcolor{red}{Ошибка:} неверное вычисление значения процентного прироста при правильном выборе исходной величины.

\point{Где именно возникает ошибка}
\textbf{Последний правильный шаг:} разность нового и старого значений делится на старое значение:
\[
\frac{600-480}{480}\cdot100\%.
\]
\textbf{Первый неверный шаг:} приравнивание этой дроби к $60000\%$. Разность в числителе равна 120, и после деления на 480 получается $0{,}25$.

\point{Почему этот шаг неверен}
Процент означает сотую часть целого. Сначала надо выяснить, какую \emph{долю от первоначальных 480 кг} составили дополнительные 120 кг, и лишь затем перевести долю в проценты. Умножение на 100 выполняется \emph{после деления}, а знак процента сохраняется в записи ответа. Значение $60000\%$ означало бы увеличение в 600 раз от исходного сбора: такой результат противоречит условию, поскольку 600 кг лишь немного больше 480 кг.

\point{Ключевая идея}
\textcolor{blue!60!black}{Ключевая идея:} процентное увеличение вычисляется по формуле
\[
\frac{\text{новое значение}-\text{старое значение}}
{\text{старое значение}}\cdot100\%.
\]

\point{Исправление от последнего правильного шага}
\[
\frac{600-480}{480}\cdot100\%
=\frac{120}{480}\cdot100\%
=\frac14\cdot100\%=25\%.
\]

\clearpage
\point{Полное правильное решение}
\begin{enumerate}
\item Находим, на сколько килограммов увеличился сбор: $600-480=120$ кг.
\item Сравниваем прибавку с количеством в прошлом году: $120:480=\frac14=0{,}25$.
\item Переводим долю в проценты: $0{,}25\cdot100\%=25\%$.
\end{enumerate}

\point{Проверка}
$25\%$ от 480 кг составляет $480\cdot0{,}25=120$ кг. Тогда $480+120=600$ кг, что совпадает с условием.

\point{Итог}
\textcolor{teal!60!black}{Итог:} сбор макулатуры увеличился на \textbf{$25\%$}.

\point{Задача на закрепление}
В прошлом месяце школьники собрали 360 кг макулатуры, в этом~--- 450 кг. На сколько процентов увеличился сбор?
\clearpage

\phantomsection\label{task:10}
\section*{Задание № 10. Проценты одной группы от другой}
\point{Текст задания}
В спортивной школе $27\%$ занимаются футболом, $44\%$~--- баскетболом, остальные~--- только неигровыми видами спорта. Сколько процентов составляет количество занимающихся неигровыми видами спорта от количества занимающихся игровыми? Ответ округлите до десятых.

\point{Ваше решение и ответ}
На представленных фотографиях решение задания № 10 и ответ отсутствуют.

\point{Ошибка}
\textcolor{red}{Ошибка:} задание оставлено без решения. Для получения ответа требуется найти доли обеих групп и сравнить их между собой.

\point{Где именно возникает неполнота}
\textbf{Что известно из условия:} футболом занимаются $27\%$, баскетболом~--- $44\%$ от всех учащихся.

\textbf{Первый недостающий шаг:} определить общую долю занимающихся игровыми видами спорта: $27\%+44\%=71\%$. Только после этого можно вычислить долю занимающихся неигровыми видами спорта.

\point{Почему важно правильно выбрать основание сравнения}
В вопросе сказано: «от количества занимающихся игровыми». Следовательно, за $100\%$ при \emph{сравнении групп} нужно принять количество занимающихся игровыми видами спорта. Проценты $27\%$, $44\%$ и оставшаяся доля первоначально даны относительно \emph{всех} учащихся. Сначала эти доли можно найти в общей шкале, а затем составить отношение нужных двух групп. Деление доли неигровых видов на $100\%$ дало бы её часть от всей школы и ответило бы на другой вопрос.

\point{Ключевая идея}
\textcolor{blue!60!black}{Ключевая идея:} «сколько процентов составляет $A$ от $B$?» означает
\[
\frac{A}{B}\cdot100\%.
\]
Делителем является именно величина, стоящая после слова «от».

\clearpage
\point{Исправление и полное правильное решение}
\begin{enumerate}
\item Доля занимающихся игровыми видами спорта:
\[
27\%+44\%=71\%.
\]
\item Доля занимающихся неигровыми видами спорта:
\[
100\%-71\%=29\%.
\]
\item Доля неигровых видов \emph{относительно} игровых:
\[
\frac{29}{71}\cdot100\%=\frac{2900}{71}\%\approx40{,}845\ldots\%.
\]
\item Округляем до одной цифры после запятой: $40{,}8\%$.
\end{enumerate}

\point{Проверка}
$40{,}8\%$ от $71\%$ всех учащихся составляет примерно $0{,}408\cdot71\%\approx28{,}97\%$, что согласуется с найденной долей $29\%$. Небольшое расхождение обусловлено округлением.

\point{Итог}
\textcolor{teal!60!black}{Итог:} количество занимающихся неигровыми видами спорта составляет \textbf{$40{,}8\%$} от количества занимающихся игровыми.

\point{Задача на закрепление}
В спортивной школе $24\%$ занимаются волейболом, $36\%$~--- баскетболом, а остальные~--- только неигровыми видами спорта. Сколько процентов составляет количество занимающихся неигровыми видами спорта от количества занимающихся игровыми? Ответ округлите до десятых.
\end{document}
